% Part: turing-machines
% Chapter: machines-computations
% Section: turing-machines

\documentclass[../../../include/open-logic-section]{subfiles}

\begin{document}

\olfileid[fr]{tur}{mac}{tur}
\olsection{Machines de Turing}

\begin{explain}
La définition formelle d'une machine de Turing
peut sembler abstraite; elle est pourtant simple.
Elle rassemble en une seule structure mathématique
toutes les données nécessaires à la description
du fonctionnement de la machine : (1) les états
dans lesquels elle peut se trouver, (2) les
symboles qui peuvent figurer sur le ruban,
(3) l'état dans lequel elle démarre et
(4) son ensemble d'instructions.
\end{explain}

\begin{defn}[Machine de Turing]
Une \emph{machine de Turing}~$M$ est un quadruplet
$\langle Q,\Sigma,q_0,\delta\rangle$ formé de :
\begin{enumerate}
\item un ensemble fini d'\emph{états}~$Q$;
\item un \emph{alphabet} fini~$\Sigma$ qui contient
  $\TMendtape$ et~$\TMblank$;
\item un \emph{état initial}~$q_0\in Q$;
\item un \emph{ensemble fini d'instructions}
  $\delta\colon Q\times\Sigma
  \pto Q\times\Sigma\times
  \{\TMleft,\TMright,\TMstay\}$.
\end{enumerate}
La fonction partielle~$\delta$ s'appelle aussi la
\emph{fonction de transition} de~$M$.
\end{defn}

\begin{explain}
Nous supposons que le ruban est infini dans une
seule direction. Il est donc utile de réserver
un symbole spécial~$\TMendtape$ pour marquer son
extrémité gauche : le programme peut ainsi
détecter qu'il risque de sortir du ruban. On
pourrait imposer que ce symbole ne soit jamais
écrasé, c'est-à-dire que, si
$\delta(q,\TMendtape)$ est définie, alors
$\delta(q,\TMendtape)=\tuple{q',\TMendtape,x}$.
Certains manuels font cette hypothèse; nous ne
la faisons pas. Pour construire une machine qui
préserve le marqueur~$\TMendtape$, il suffit de
prendre les précautions nécessaires. Autoriser
son écrasement offre aussi, dans certains cas,
une souplesse utile.
\end{explain}

\begin{ex}
\emph{Machine paire} : formellement, cette machine
est le quadruplet $\tuple{Q,\Sigma,q_0,\delta}$,
où
\begin{align*}
Q & = \{ q_0, q_1 \} \\
\Sigma & = \{ \TMendtape, \TMblank, \TMstroke \}, \\
\delta(q_0, \TMstroke) & = \tuple{q_1, \TMstroke, \TMright},\\
\delta(q_1, \TMstroke) & = \tuple{q_0, \TMstroke, \TMright},\\
\delta(q_1, \TMblank)  & = \tuple{q_1, \TMblank, \TMright}.
\end{align*}
\end{ex}

\end{document}
